% Promoted from experimental/experiments.tex lines 1944--2090.
% Status: stable high-agreement note.  This is a TeX fragment intended to be
% included from experimental/experiments.tex or another paper-level wrapper.
% Dependencies: theorem environments, standard macros such as \F and \RS,
% and earlier local labels referenced inside the note.

\subsection{The \(\F_{17^{32}}\), \(n=512\), \(k=256\) protocol-facing ledger}

\begin{corollary}[high-agreement adjacent ledgers for the \(17^{32}\) row]
\label{cor:f17-adjacent-ledgers}
Let \(H\subset\F_{17^{32}}^*\) have order \(512\), and let
\[
  C=\RS[\F_{17^{32}},H,256].
\]
Put \(q=17^{32}\) and
\[
  B_*=\left\lfloor q/2^{128}\right\rfloor=6.
\]
Then the following statements hold.

\begin{enumerate}
\item For every \(a\ge427\),
\[
  \operatorname{LD}_{\rm ca}(C,a)
  =
  \operatorname{LD}_{\rm sw,proj}(C,a)
  =
  \operatorname{LD}_{\rm sw}(C,a)
  =
  513-a.
\]
Thus the affine/projective line numerator is \(6\) at \(a=507\) and \(7\) at
\(a=506\).

\item For every \(\mu\ge1\) and every \(a\ge384\),
\[
  \Lambda_\mu(C,a)=1.
\]

\item For degree-\(d\) finite-parameter curves, whenever
\[
  a\ge
  \left\lceil\frac{(d+1)512+256}{d+2}\right\rceil,
\]
we have
\[
  \operatorname{CurveLD}_{\rm sw}^{(d)}(C,a)
  =
  \operatorname{CurveLD}_{\rm ca}^{(d)}(C,a)
  =
  \min\{17^{32},\,d(513-a)\}.
\]
\end{enumerate}
\end{corollary}

\begin{proof}
The first item is Corollary~\ref{cor:f17-32-exact-tangent-gate} together with
Theorem~\ref{thm:ca-projective-tangent-staircase}.  The second item is
Theorem~\ref{thm:interleaved-unique-high-agreement}, because
\(2a-512\ge256\) is equivalent to \(a\ge384\).  The third item is
Theorem~\ref{thm:curve-tangent-staircase} after substituting \(n=512\) and
\(k=256\).  The integer \(B_*=6\) follows from
\[
  6\cdot2^{128}<17^{32}<7\cdot2^{128}.
\]
\end{proof}

\begin{corollary}[conditional high-agreement protocol ledger]
\label{cor:high-agreement-protocol-ledger}
Consider a protocol reduction for the row of
Corollary~\ref{cor:f17-adjacent-ledgers} whose coding-theoretic error at
agreement \(a\) is exactly
\[
  \frac{N_{\rm geom}(a)}{17^{32}}
  +
  \frac{\Lambda_\mu(C,a)}{17^{32}}
  +
  \varepsilon_{\rm query},
\]
where \(N_{\rm geom}(a)\) is one affine/projective line CA/MCA numerator or one
degree-\(d\) finite-parameter curve numerator, and where no additional fold,
hash, Fiat--Shamir, or extension-lift loss is hidden in the notation.  In the
high-agreement ranges above, this ledger becomes the following explicit integer
test.

For affine/projective line CA/MCA,
\[
  N_{\rm geom}(a)+\Lambda_\mu(C,a)
  =
  (513-a)+1
  =
  514-a
  \qquad (a\ge427).
\]
With \(\varepsilon_{\rm query}=0\), this is at most the \(2^{-128}\) budget iff
\(a\ge508\).  Thus the combined line-plus-list ledger is safe at agreement
\(508\) and unsafe at agreement \(507\).

For degree-\(d\) finite-parameter curve MCA/CA,
\[
  N_{\rm geom}(a)+\Lambda_\mu(C,a)
  =
  d(513-a)+1
\]
in the curve-staircase range.  With \(\varepsilon_{\rm query}=0\), the safe
grid points are exactly those satisfying
\[
  d(513-a)+1\le6 .
\]
In particular:
\[
\begin{array}{c|c|c}
d & \text{curve-plus-list safe grid} & \text{curve term alone safe grid}\\
\hline
1 & a\ge508\ (r\le4) & a\ge507\ (r\le5)\\
2 & a\ge511\ (r\le1) & a\ge510\ (r\le2)\\
3 & a=512\ (r=0) & a\ge511\ (r\le1)\\
4 & a=512\ (r=0) & a=512\ (r=0)\\
5 & a=512\ (r=0) & a=512\ (r=0)\\
6 & \text{no safe grid point with the list term} & a=512\ (r=0)\\
d\ge7 & \text{no safe grid point with the list term} &
        \text{no safe grid point for the curve term alone}
\end{array}
\]
where \(r=n-a\) is the integer Hamming radius.
\end{corollary}

\begin{proof}
By Corollary~\ref{cor:f17-adjacent-ledgers}, the interleaved-list numerator is
\(1\) throughout the displayed high-agreement range.  The affine/projective line
numerator is \(513-a\), and the degree-\(d\) curve numerator is \(d(513-a)\).
Since
\[
  \left\lfloor 17^{32}/2^{128}\right\rfloor=6,
\]
the code-level numerator must be at most \(6\) when
\(\varepsilon_{\rm query}=0\).  The displayed inequalities are exactly the
resulting integer conditions.
\end{proof}

\begin{remark}[what this does and does not settle]
Theorems~\ref{thm:ca-projective-tangent-staircase},
\ref{thm:curve-tangent-staircase}, and
\ref{thm:interleaved-unique-high-agreement} are unconditional finite
statements under the hypotheses printed above: the CA/projective/curve
statements are Reed--Solomon evaluation-domain statements, while interleaved
uniqueness is MDS.  Corollary~\ref{cor:high-agreement-protocol-ledger} is
conditional on the protocol really consuming only the printed coding terms over
the printed field.  It does not cover reductions with additional curve
families, extension-line lifts, folding losses, query errors, or cryptographic
errors unless those terms are added to the ledger.
\end{remark}
